% =====================================================================
%  HOW TO USE
%  1. Go to overleaf.com -> New Project -> Templates -> search "ACL 2023"
%  2. Open as template, then REPLACE the template's main.tex with this file.
%  3. Also copy references.bib into the project.
%  4. Recompile.
%  (This file assumes the ACL style files acl.sty / acl_natbib.bst exist,
%   which they do in the ACL template project.)
% =====================================================================

\documentclass[11pt]{article}
\usepackage[review]{acl}      % use [final] when submitting
\usepackage{times}
\usepackage{latexsym}
\usepackage{booktabs}
\usepackage{graphicx}
\usepackage{amsmath}
\usepackage{multirow}
\usepackage[T1]{fontenc}
\usepackage[utf8]{inputenc}
\usepackage{microtype}

\title{Difficulty-Aware Language Routing for Korean \\ Mathematical Reasoning in Small Language Models}

% Fill in your names / student IDs / course
\author{Your Name \\
  Department / University \\
  \texttt{email@domain} \\\And
  Teammate Name \\
  Department / University \\
  \texttt{email@domain}}

\begin{document}
\maketitle

% =====================================================================
\begin{abstract}
% DRAFT --- write this LAST, after everything else is done.
Small language models reason markedly better in English than in Korean
on mathematical problems, in part because even strong teacher models
produce lower-quality chain-of-thought (CoT) in Korean than in English.
We investigate whether a model's stronger English reasoning can be
leveraged to improve its Korean mathematical reasoning. We introduce
\textbf{Difficulty-Aware Language Routing (DALR)}, a data-construction
method that selects, per problem, whether to learn from a Korean or
English teacher CoT based on teacher correctness. We further combine
specialist models via \emph{model soup} \citep{wortsman2022soup} and a
\emph{cross-lingual ensemble} at inference. On HRM8K (Korean) and GSM8K
(English), DALR improves Korean accuracy over a Korean-only baseline,
and an ablation shows the gain stems from \emph{where} English CoT is
routed (hard problems) rather than from simply adding more data.
% TODO: final numbers once 1,319 eval done.
\end{abstract}

% =====================================================================
\section{Introduction}
\label{sec:intro}

Large language models have achieved strong mathematical reasoning, but
small models (e.g., 3B parameters) lag substantially, especially in
lower-resource languages. We observe a large gap between English and
Korean math accuracy for a 3B model: the base Qwen2.5-3B-Instruct
solves \textbf{67.4\%} of English (GSM8K) but only \textbf{52.4\%} of
Korean (HRM8K) problems.

A natural fix is knowledge distillation: have a strong teacher generate
chain-of-thought (CoT) solutions and fine-tune the small student. But
this faces a subtle obstacle: \emph{the teacher itself is weaker in
Korean}. In our data, Gemini produces correct CoT for \textbf{97.5\%}
of English problems but only \textbf{85.7\%} of Korean problems. Naively
distilling Korean CoT therefore caps the student at a degraded signal,
while English keeps improving --- widening the gap.

This motivates our central question:
\begin{quote}
\emph{Can a model's stronger English reasoning be used to improve its
Korean mathematical reasoning?}
\end{quote}
Our key assumption is that mathematical \emph{logic} is largely
language-agnostic: the steps to solve a problem are the same whether it
is stated in Korean or English. If English reasoning patterns can be
connected to Korean problems, Korean performance should improve.

We approach this through \emph{quality-aware aggregation} of
cross-lingual signals at three levels:
\begin{itemize}
  \item \textbf{Data level (DALR).} Per problem, route to a Korean or
  English teacher CoT based on teacher correctness
  (Section~\ref{sec:dalr}).
  \item \textbf{Model level (soup).} Average the weights of specialist
  models \citep{wortsman2022soup} (Section~\ref{sec:soup}).
  \item \textbf{Inference level (ensemble).} Majority-vote across models
  reasoning in different languages (Section~\ref{sec:clsc}).
\end{itemize}

Our contributions:
\begin{enumerate}
  \item We propose \textbf{DALR}, a difficulty-aware data-construction
  method for cross-lingual CoT distillation, and show via a controlled
  ablation that its gains come from \emph{routing}, not data scaling.
  \item We apply and analyze model soup and cross-lingual ensembling in
  the Korean math setting, characterizing a Korean--English trade-off
  and how aggregation resolves it.
  \item We provide a statistical analysis (bootstrap CIs, McNemar tests)
  and a per-difficulty breakdown of where each method helps.
\end{enumerate}

% =====================================================================
\section{Related Work}
\label{sec:related}

\paragraph{CoT distillation.} % TODO cite a couple
Chain-of-thought prompting \citep{wei2022cot} and its distillation into
smaller models...

\paragraph{Cross-lingual reasoning.} % TODO cite xCoT etc.
Prior work transfers reasoning across languages via translation or
cross-lingual instruction tuning. Unlike methods that augment prompts,
DALR selects training CoT \emph{language} per problem by teacher quality.

\paragraph{Weight averaging \& self-consistency.}
Model soup \citep{wortsman2022soup} averages fine-tuned weights for a
single low-cost model. Self-consistency \citep{wang2022selfconsistency}
samples multiple chains from one model. We instead ensemble distinct
models trained in different languages.

% =====================================================================
\section{Method}
\label{sec:method}

\subsection{Setup and Data}
\label{sec:data}
We use Qwen2.5-3B-Instruct as the student and Gemini as the teacher.
Training problems come from GSM8K (7{,}473) and its Korean translation
GSM8K-ko. We generate teacher CoT in both languages and keep only
correct ones (English: 7{,}283; Korean: 6{,}407). We evaluate on HRM8K
(human-reviewed Korean, 1{,}319) and GSM8K (English, 1{,}319).
All students are trained with LoRA ($r{=}32$).

\subsection{Difficulty-Aware Language Routing (DALR)}
\label{sec:dalr}
For each Korean training problem, we route based on teacher correctness:
\begin{itemize}
  \item Korean teacher correct $\rightarrow$ learn from \textbf{Korean CoT}.
  \item Korean wrong, English correct $\rightarrow$ learn from
  \textbf{English CoT} (a \emph{bridge}: Korean question, English solution).
  \item Both wrong $\rightarrow$ \textbf{discard} (noise).
\end{itemize}
This yields 6{,}407 Korean and 920 English-bridge examples (7{,}327
total). Intuitively, teacher failure in Korean signals that the problem
is hard to reason about in Korean, so a reliable English solution is a
better learning target.

\subsection{Model Soup}
\label{sec:soup}
Different training recipes produce \emph{specialists}: English-strong,
Korean-strong, or balanced. We combine them by element-wise averaging of
their LoRA weights \citep{wortsman2022soup}, producing one model at
$1\times$ inference cost.

\subsection{Cross-Lingual Ensemble (CLSC)}
\label{sec:clsc}
At inference we run multiple models on the same problem and take a
majority vote, breaking ties by a designated model. Because models
reason in different languages, their errors differ, improving ensemble
robustness.

% =====================================================================
\section{Experiments}
\label{sec:exp}
We evaluate single greedy decoding with a zero-shot CoT system prompt
(repetition penalty $1.1$). Our setups are: A (base), B (English CoT),
C (Korean CoT), D (bilingual mix), E (two-stage), F (DALR), and
F\_random (DALR ablation with randomly placed English bridges).

% =====================================================================
\section{Results}
\label{sec:results}

\begin{table}[t]
\centering
\small
\begin{tabular}{lcc}
\toprule
\textbf{Setup} & \textbf{HRM8K (KO)} & \textbf{GSM8K (EN)} \\
\midrule
A (Base)        & 53.9  & TBD   \\
B (EN CoT)      & 54.4  & TBD   \\
C (KO CoT)      & 59.5  & 70.7  \\
D (Bilingual)   & TBD   & TBD   \\
E (Two-stage)   & TBD   & TBD   \\
F (DALR)        & \textbf{61.8} & 69.4 \\
F\_random       & 58.8  & 72.7  \\
\midrule
soup (B+C+D+E+F) & 60.1 & \textbf{78.6} \\
\midrule
CLSC 6-way       & 68.6$^\dagger$ & 81.6$^\dagger$ \\
CLSC 7-way (+soup) & \textbf{70.0}$^\dagger$ & \textbf{84.4}$^\dagger$ \\
\bottomrule
\end{tabular}
\caption{Accuracy (\%) on the full 1{,}319 test sets.
$^\dagger$ preliminary (500-sample) numbers; full-set update pending.
TBD entries are running.}
\label{tab:main}
\end{table}

\paragraph{DALR ablation (routing vs. data scaling).}
F and F\_random use the \emph{same amount} of English CoT (920); they
differ only in \emph{where} it is placed. F routes English to hard
(Korean-failed) problems; F\_random places it on random easy problems.
On HRM8K, F (61.8) $>$ F\_random (58.8) by $+3.0$ points (McNemar
$p{=}0.030$), and F\_random even falls below the Korean-only baseline C
(59.5). This shows the improvement comes from \emph{routing}, not from
adding English data.

\paragraph{Trade-off and its resolution.}
DALR improves Korean but slightly lowers English (F: 69.4 vs.\ C: 70.7).
Model soup recovers English strongly (78.6, $p{<}0.001$ vs.\ F) while
keeping competitive Korean, and the cross-lingual ensemble achieves the
best overall results.

% =====================================================================
\section{Analysis}
\label{sec:analysis}

\paragraph{Per-difficulty breakdown.}
Stratifying HRM8K problems by how many setups solve them, F outperforms
the soup on \emph{medium-difficulty} problems % TODO numbers
(those solved by 2--3 of 4 models), suggesting DALR's English bridge
helps precisely where Korean-only reasoning is unreliable.

\paragraph{CoT language at inference.} % optional
% TODO: language usage table if you want it.

% =====================================================================
\section{Conclusion}
\label{sec:conclusion}
We presented DALR, a difficulty-aware routing method for cross-lingual
CoT distillation, and showed---through a controlled ablation---that its
benefit stems from routing English reasoning to hard Korean problems.
Combined with model soup and cross-lingual ensembling, our approach
improves small-model Korean mathematical reasoning.

\section*{Limitations}
HRM8K is a human-reviewed translation of GSM8K rather than natively
Korean math, so gains may partly reflect translated-problem solving.
The DALR-vs-Korean-only gap is modest and not always significant at our
sample size. We use a single teacher and a single 3B student.

% =====================================================================
\bibliography{references}
\bibliographystyle{acl_natbib}

\end{document}
